\documentclass[11pt]{article}

% -------------------- Packages --------------------
\usepackage[a4paper,margin=1in]{geometry}
\usepackage{hyperref}
\usepackage{enumitem}
\usepackage{titlesec}
\usepackage{fancyhdr}
\usepackage{xcolor}
\usepackage{listings}
% listings has no built-in yaml or Docker language; define them (plain, no highlighting) so latexmk builds without errors.
\lstdefinelanguage{yaml}{}
\lstdefinelanguage{Docker}{}
\usepackage{courier}

% -------------------- Hyperref Setup --------------------
\hypersetup{
  colorlinks=true,
  linkcolor=blue,
  urlcolor=blue,
  citecolor=blue
}

% -------------------- Header/Footer --------------------
\pagestyle{fancy}
\fancyhf{}
\lhead{MCV Lab}
\rhead{CI/CD, Containers, and MLOps}
\cfoot{\thepage}

% -------------------- Section Formatting --------------------
\titleformat{\section}{\Large\bfseries}{\thesection}{0.8em}{}
\titleformat{\subsection}{\normalsize\bfseries}{\thesubsection}{0.8em}{}
\titleformat{\subsubsection}{\normalsize\bfseries}{\thesubsubsection}{0.8em}{}

\titlespacing*{\section}{0pt}{1.5em}{0.8em}
\titlespacing*{\subsection}{0pt}{1.0em}{0.5em}
\titlespacing*{\subsubsection}{0pt}{0.8em}{0.4em}

% -------------------- Global List Formatting --------------------
\setlist[itemize]{
  label={-},
  leftmargin=2em,
  labelsep=0.6em,
  itemsep=0.2em
}
\setlist[itemize,2]{itemsep=0.15em}
\setlist[enumerate]{leftmargin=2em,itemsep=0.15em}

% -------------------- Listings / Code Style --------------------
\definecolor{codegray}{rgb}{0.95,0.95,0.95}
\definecolor{commentgreen}{rgb}{0.0,0.5,0.0}

\lstset{
  basicstyle=\footnotesize\ttfamily,
  backgroundcolor=\color{codegray},
  frame=single,
  breaklines=true,
  showstringspaces=false,
  keywordstyle=\color{blue},
  commentstyle=\color{commentgreen},
  stringstyle=\color{red},
  columns=fullflexible,
  keepspaces=true,
  tabsize=2,
  xleftmargin=0.02\textwidth,
  xrightmargin=0.02\textwidth
}

% -------------------- Document Metadata --------------------
\title{\textbf{Medical Computer Vision Lab \\ \small{
CI/CD, Containerization, and MLOps Practices: Lab Standard}}}
\author{Maintainer: Ashery Mbilinyi}
\date{Version 1.0 : \today}

\begin{document}
\maketitle

\tableofcontents
\newpage

\section{Purpose}
This document defines how we operationalize machine learning (ML) systems in the MCV Lab by combining:
\begin{itemize}
  \item \textbf{CI/CD} (automation of testing, integration, and delivery)
  \item \textbf{Containerization} (Docker, Kubernetes) for portability and scalability
  \item \textbf{MLOps practices} for full lifecycle model management
\end{itemize}

The goal is to ensure our ML work is:
\begin{itemize}
  \item \textbf{Reproducible}: experiments can be rerun months later with identical setup
  \item \textbf{Reliable}: training and inference are stable and testable
  \item \textbf{Deployable}: models can move from research to production systematically
  \item \textbf{Maintainable}: changes are reviewed, tested, and traceable
\end{itemize}

\section{Key Concept Summary (One Page)}
\subsection{Core Definitions}
\begin{itemize}
  \item \textbf{CI (Continuous Integration)}: automatically test and validate code changes on every push/PR.
  \item \textbf{CD (Continuous Delivery/Deployment)}: automatically package and deliver validated code to environments (staging/prod).
  \item \textbf{Containerization}: package the runtime environment so code runs consistently across machines.
  \item \textbf{Docker}: toolchain to build/run containers locally or on servers.
  \item \textbf{Kubernetes (K8s)}: orchestration platform to run containers reliably across a cluster.
  \item \textbf{MLOps}: practices/tools for managing the end-to-end lifecycle of ML models (data $\rightarrow$ training $\rightarrow$ evaluation $\rightarrow$ deployment $\rightarrow$ monitoring).
\end{itemize}

\subsection{The Relationship (Mental Model)}
\begin{itemize}
  \item \textbf{Docker} answers: \emph{``How do we package and run this ML system reproducibly?''}
  \item \textbf{Kubernetes} answers: \emph{``How do we run this at scale and keep it alive reliably?''}
  \item \textbf{CI/CD} answers: \emph{``How do we automate testing and delivery when code changes?''}
  \item \textbf{MLOps} answers: \emph{``How do we manage models like real software products?''}
\end{itemize}

\section{Containerization in ML}

\subsection{What Containerization Solves}
ML projects often break due to environment mismatches:
\begin{itemize}
  \item different CUDA versions
  \item different PyTorch versions
  \item missing system libraries (OpenCV, libjpeg, etc.)
  \item dependency conflicts across projects
\end{itemize}

A container solves this by packaging:
\begin{itemize}
  \item OS-level dependencies
  \item Python dependencies
  \item runtime settings and entrypoint scripts
\end{itemize}

\subsection{Docker: Packaging and Reproducible Execution}
\textbf{Docker} is used to:
\begin{itemize}
  \item build container images from a \texttt{Dockerfile}
  \item run training/evaluation in isolated environments
  \item support consistent onboarding of new lab members
\end{itemize}

\subsubsection{Minimum Docker Standard (Lab Policy)}
Every ML project repository should include:
\begin{itemize}
  \item a \texttt{Dockerfile}
  \item pinned dependencies (\texttt{requirements.txt}, \texttt{poetry.lock}, or \texttt{environment.yml})
  \item documented run commands in \texttt{README.md}
\end{itemize}

\subsection{Kubernetes: Orchestration for Scale}
\textbf{Kubernetes (K8s)} manages running containers across a cluster.
It enables:
\begin{itemize}
  \item scheduling jobs across multiple nodes
  \item automatic restart/reschedule if a node fails
  \item autoscaling (e.g., increase number of inference replicas)
  \item rolling updates (deploy new version with zero downtime)
\end{itemize}

\subsubsection{When Kubernetes is Relevant for the Lab}
Kubernetes is most useful when:
\begin{itemize}
  \item running a shared GPU cluster platform
  \item serving models as APIs used by others
  \item running continuous pipelines with multiple services (training + registry + evaluation + serving)
\end{itemize}

\subsubsection{When Kubernetes is Not Necessary}
Kubernetes is often \textbf{not required} when:
\begin{itemize}
  \item training is run on single machines or HPC nodes
  \item the lab uses SLURM scheduling (academic default)
  \item no persistent production model serving is needed
\end{itemize}

\section{CI/CD in ML (What It Means in Practice)}

\subsection{CI: Continuous Integration}
CI is an automation pipeline triggered by:
\begin{itemize}
  \item every commit / push
  \item every pull request (PR)
\end{itemize}

CI ensures that the codebase stays healthy and deployable.

\subsubsection{Typical CI Checks in ML Projects}
\begin{itemize}
  \item \textbf{Code quality}: linting (ruff/flake8), formatting (black), type-checking (mypy)
  \item \textbf{Unit tests}: dataset loaders, transforms, model forward pass
  \item \textbf{Training sanity tests}: run 5--20 iterations on tiny data and confirm:
    \begin{itemize}
      \item no runtime crash
      \item loss decreases
      \item outputs have expected shapes
    \end{itemize}
  \item \textbf{Reproducibility tests} (optional but recommended):
    \begin{itemize}
      \item fixed seed yields consistent results for small runs
    \end{itemize}
\end{itemize}

\subsection{CD: Continuous Delivery / Deployment}
CD runs after CI passes. It automates:
\begin{itemize}
  \item building Docker images
  \item pushing images to registries (GHCR, DockerHub, GCP Artifact Registry)
  \item deploying to staging
  \item optionally deploying to production
\end{itemize}

\subsubsection{CD in ML Has Two Meanings}
In ML, delivery can refer to:
\begin{itemize}
  \item \textbf{Code delivery}: deploy training/inference software
  \item \textbf{Model delivery}: promote and deploy a specific trained model artifact
\end{itemize}

\subsubsection{Safe Deployment Pattern}
\begin{itemize}
  \item deploy new version to staging
  \item run smoke tests + basic inference tests
  \item deploy gradually (canary rollout) when possible
  \item keep rollback path (\texttt{model\_vX} $\rightarrow$ \texttt{model\_v(X-1)})
\end{itemize}

\section{MLOps: Full Model Lifecycle Management}

\subsection{Why MLOps Is Needed}
ML differs from classic software:
\begin{itemize}
  \item models depend on data quality and distributions
  \item performance can degrade over time (data drift)
  \item ``bugs'' may appear only as metric drops
  \item retraining is often required
\end{itemize}

MLOps manages ML as a \textbf{continuous lifecycle}, not a one-time training script.

\subsection{Core Components of MLOps}
\subsubsection{(1) Data Management}
\begin{itemize}
  \item version datasets and splits (train/val/test)
  \item document preprocessing and transforms
  \item prevent leakage across splits
  \item record dataset metadata and provenance
\end{itemize}

\subsubsection{(2) Experiment Tracking}
Every run must log:
\begin{itemize}
  \item configuration (hyperparameters, seed)
  \item metrics (loss, AUROC, F1, etc.)
  \item code version (git commit hash)
  \item artifacts (plots, confusion matrix)
\end{itemize}

Recommended tools:
\begin{itemize}
  \item TensorBoard (baseline)
  \item Weights \& Biases (W\&B)
  \item MLflow
\end{itemize}

\subsubsection{(3) Model Checkpointing and Model Registry}
We distinguish:
\begin{itemize}
  \item \textbf{Checkpoints}: intermediate training states (\texttt{last.ckpt}, \texttt{epoch\_10.ckpt})
  \item \textbf{Model version}: approved/published model artifact with metadata
\end{itemize}

A model registry tracks:
\begin{itemize}
  \item which model is in production
  \item which dataset/training config produced it
  \item evaluation metrics and approval status
\end{itemize}

\subsubsection{(4) Automated Pipelines}
Automate training + evaluation:
\begin{itemize}
  \item ``train'' stage
  \item ``evaluate'' stage
  \item ``promote model'' stage if metrics pass
\end{itemize}

Typical tools:
\begin{itemize}
  \item Airflow, Kubeflow Pipelines, Vertex AI Pipelines
\end{itemize}

\subsubsection{(5) Deployment and Serving}
Deploying a model means running inference:
\begin{itemize}
  \item batch inference jobs (HPC or cloud)
  \item online inference (REST/gRPC service)
  \item integration into clinical/research pipelines
\end{itemize}

Typical deployment tools:
\begin{itemize}
  \item FastAPI + Docker (lab internal)
  \item BentoML / TorchServe / Triton (advanced)
  \item KServe / Seldon (K8s-native)
\end{itemize}

\subsubsection{(6) Monitoring}
Production monitoring includes:
\begin{itemize}
  \item latency (response time)
  \item error rates
  \item input data drift
  \item prediction drift
  \item degradation in downstream clinical performance (when labels exist)
\end{itemize}

\section{End-to-End Workflow (How Everything Connects)}

\subsection{Standard Flow}
\begin{enumerate}
  \item Researcher proposes change (new model / preprocessing / metric).
  \item Code is pushed as a PR.
  \item \textbf{CI runs}: unit tests + sanity training + linting.
  \item If CI passes: \textbf{CD builds} a Docker image and publishes it.
  \item Image is deployed to staging (optional in academic setting).
  \item Training jobs run using:
    \begin{itemize}
      \item Docker/Apptainer image (environment)
      \item SLURM/K8s (execution scheduling)
    \end{itemize}
  \item Outputs logged to tracking system (TensorBoard/W\&B/MLflow).
  \item Best model is promoted into registry and deployed.
  \item Monitoring detects drift and triggers retraining when needed.
\end{enumerate}

\subsection{What This Achieves}
\begin{itemize}
  \item fewer broken experiments
  \item fewer wasted GPU hours
  \item faster onboarding
  \item traceable results for publications
  \item reliable path from research to deployment
\end{itemize}

\section{Lab Policy Checklist}

\subsection{Minimum Requirements (All Projects)}
\begin{itemize}
  \item Training runs without code edits (CLI + YAML/JSON configs)
  \item Logs saved per run (metrics + config + command)
  \item Checkpoints saved (\texttt{last} + \texttt{best})
  \item Deterministic seed option exists
  \item Container recipe exists (\texttt{Dockerfile})
\end{itemize}

\subsection{Recommended Requirements (Research at Scale)}
\begin{itemize}
  \item CI pipeline (GitHub Actions recommended)
  \item Automated training sanity tests in CI
  \item Experiment tracking (W\&B or MLflow)
  \item Model registry practice (artifact versioning)
\end{itemize}

\subsection{Optional (When Running Services)}
\begin{itemize}
  \item staging environment for deployments
  \item monitoring for drift/performance
  \item Kubernetes-based serving when multi-tenant
\end{itemize}

\section{Appendix: Minimal Templates}

\subsection{Example: Minimal GitHub Actions CI}
Create: \texttt{.github/workflows/ci.yml}
\begin{lstlisting}[language=yaml]
name: CI

on:
  push:
  pull_request:

jobs:
  test:
    runs-on: ubuntu-latest
    steps:
      - uses: actions/checkout@v4

      - name: Set up Python
        uses: actions/setup-python@v5
        with:
          python-version: "3.10"

      - name: Install dependencies
        run: |
          pip install -r requirements.txt

      - name: Lint
        run: |
          python -m compileall .

      - name: Unit tests
        run: |
          pytest -q
\end{lstlisting}

\subsection{Example: CD Builds and Pushes Docker Image}
\begin{lstlisting}[language=yaml]
name: CD

on:
  push:
    branches: [ "main" ]

jobs:
  build-and-push:
    runs-on: ubuntu-latest
    steps:
      - uses: actions/checkout@v4
      - name: Build Docker image
        run: docker build -t mcv-project:latest .
\end{lstlisting}

\subsection{Example: Minimal Dockerfile for ML}
\begin{lstlisting}[language=Docker]
FROM pytorch/pytorch:2.3.1-cuda12.1-cudnn8-runtime

WORKDIR /workspace

RUN apt-get update && apt-get install -y --no-install-recommends \
    git curl ca-certificates \
    && rm -rf /var/lib/apt/lists/*

COPY requirements.txt /workspace/requirements.txt
RUN pip install --no-cache-dir -r requirements.txt

COPY . /workspace

CMD ["python", "scripts/train.py", "--help"]
\end{lstlisting}

\end{document}
